\documentclass[11pt]{article}

\usepackage[margin=1in]{geometry}
\usepackage{amsmath, amssymb, amsthm}
\usepackage{booktabs}
\usepackage{enumitem}
\usepackage{xcolor}
\usepackage[colorlinks=true, linkcolor=blue!60!black, urlcolor=blue!60!black,
            citecolor=blue!60!black]{hyperref}
\usepackage{listings}

\lstset{
  basicstyle=\ttfamily\small,
  breaklines=true,
  columns=fullflexible,
  keepspaces=true,
  frame=single,
  framesep=4pt,
  backgroundcolor=\color{gray!7},
  xleftmargin=6pt,
}

\newcommand{\code}[1]{\texttt{#1}}
\newcommand{\bigO}{\mathcal{O}}

\title{Making \texttt{ORJointUpdate} Faster:\\
A Cost Model and Ranked Set of Speedups}
\author{Notes on \texttt{OR\_joint\_update.Rmd}}
\date{15 July 2026}

\begin{document}
\maketitle

\begin{abstract}
The joint disturbance/biomass filter \code{ORJointUpdate} is slow because it runs
the VSEM forest model once per ensemble member, per mixture component, per
Gauss--Laguerre precision node, per timestep, per site. This note builds an explicit
cost model, identifies the single dominant term, and lists concrete speedups ranked by
expected impact. The headline result is that the precision quadrature (a factor of
$G\!\approx\!20$) currently multiplies the VSEM cost even though \emph{VSEM does not
depend on the precision node} --- so most of that work is redundant and can be removed
either exactly (parallelism, pruning, batching) or with a well-justified approximation
(Rao--Blackwellizing the precision, sigma-point forecasting).
\end{abstract}

\section{Where the time goes: a cost model}
\label{sec:cost}

The expensive object is a single 365-day VSEM simulation. Every other operation
(Kalman update of a $3\times3$ system, moment-matching collapse, log-sum-exp over
$G$ nodes) is cheap by comparison. Counting VSEM calls therefore counts the cost.

The forecast step \code{forecast\_split} (\code{OR\_joint\_update.Rmd:330}) contains
the hotspot,
\begin{lstlisting}
for (j in 1:J) Nf[j, ] <- VSEM(THETA[j, ], PAR = PAR[tsel])[365, 2:4]
\end{lstlisting}
and \code{forecast\_split} is called inside three nested loops --- over timesteps
$t=1{:}NT$, precision nodes $g=1{:}G$, and input components $c=1{:}K_{\text{in}}$
(with $K_{\text{in}}=1$ at $t=1$ and $K_{\text{in}}=D$ afterwards). The whole
\code{ORJointUpdate} call is itself run once per site. Hence

\begin{equation}
\label{eq:cost}
\underbrace{N_{\text{VSEM}}}_{\text{365-day sims}}
\;\approx\;
N_{\text{sites}} \cdot NT \cdot G \cdot D \cdot J,
\end{equation}

where
\begin{itemize}[nosep]
  \item $J = \code{ne} = \code{nrow(Xic)}$ is the ensemble size per component/node;
  \item $D = \code{dist\$n}+1 = 5$ (one undisturbed label $+$ four disturbance
        classes: other, cut, fire, biotic);
  \item $G = 20$ is the number of Gauss--Laguerre nodes for the precision $\tau$;
  \item $NT \approx 33$ timesteps.
\end{itemize}

With these defaults a single site costs $\approx 33\cdot20\cdot5\cdot J = 3300\,J$
VSEM runs, and the site loop (\code{run\_joint\_update\_sites},
\code{OR\_joint\_update.Rmd:588}) runs them one site at a time. Every factor in
Eq.~\eqref{eq:cost} is a lever. The rest of this note attacks them in order of payoff.

\paragraph{Summary of expected gains.}
\begin{center}
\begin{tabular}{@{}llll@{}}
\toprule
\# & Idea & Expected speedup & Kind \\
\midrule
1 & Rao--Blackwellize the precision (drop the per-node VSEM) & $\sim G \approx 20\times$ & approximation \\
2 & Sigma-point / EKF forecast instead of $J$-member MC     & $\sim J/7$ (often $10\text{--}15\times$) & approximation \\
3 & Prune negligible nodes and disturbance labels           & $2\text{--}5\times$ (data-dependent) & near-exact \\
4 & Parallelize sites (and the $g$/$j$ loops)               & $\sim$ \#cores & exact \\
5 & Batch/Rcpp VSEM (one call for all members)              & $2\text{--}10\times$ constant & exact \\
6 & Micro-optimizations (preallocation, hoisting)           & $1.2\text{--}2\times$ & exact \\
\bottomrule
\end{tabular}
\end{center}
Ideas 1 and 2 are multiplicative with each other and with 3--6. Combining 1, 2 and 4
would turn $G\cdot D\cdot J$ VSEM runs per step into $D\cdot 7$ runs spread over all
cores --- roughly a $\bigO(10^2\text{--}10^3)$ overall reduction.

\section{Idea 1 --- Rao--Blackwellize the precision quadrature}
\label{sec:rb}

\paragraph{The observation.} VSEM's output is a deterministic function of the state
ensemble (written into \code{THETA[,9:11]}) and the driver \code{PAR[tsel]}. It does
\emph{not} depend on the precision node $\tau_g$. In the current code $\tau_g$ enters
only \emph{after} VSEM, as the analytically-added process error on the wood pool,
\begin{lstlisting}
Qg <- Q; Qg[3, 3] <- 1 / tau_g          # OR_joint_update.Rmd:334
...
cSig[[1]] <- Sig.ens + Qg               # OR_joint_update.Rmd:348
\end{lstlisting}
So of the $G\cdot D\cdot J$ VSEM runs per step, the $G$ factor is almost entirely
redundant work: the same nonlinear pushforward is recomputed 20 times and then
perturbed by 20 different additive $3\times3$ covariances.

\paragraph{Why it is currently repeated.} The filter carries a \emph{separate
ensemble per node}, \code{ens[,,c,g]} (\code{OR\_joint\_update.Rmd:371}), because the
node-specific analysis posterior --- resampled at
\code{OR\_joint\_update.Rmd:488--494} --- differs by $g$ through $Q_g$. Different
incoming ensembles $\Rightarrow$ different VSEM inputs $\Rightarrow$ VSEM must run
again. The precision has leaked into the Monte-Carlo state representation.

\paragraph{The fix.} Move $\tau$ out of the ensemble and into the analytic covariance
layer, i.e.\ Rao--Blackwellize (marginalize) the precision. Concretely:
\begin{enumerate}[nosep]
  \item Carry a \emph{single}, node-independent analysis mixture forward (per label
        $k$, per input component $c$) --- the $\tau$-marginal posterior --- rather
        than $G$ copies.
  \item Run VSEM once per $(c)$ on that shared ensemble to get the pushforward mean
        and covariance $(\hat\mu_c,\ \Sigma^{\text{ens}}_c)$.
  \item For each node $g$, form the forecast covariance
        $\Sigma^{\text{ens}}_c + Q_g$ \emph{analytically} (exactly what line 348 already
        does), run the Kalman update, and accumulate the node evidence $L_t(\tau_g)$
        and posterior weight $W_g$ as today.
\end{enumerate}
This turns the VSEM count from $G\cdot D\cdot J$ into $D\cdot J$ per step --- a factor
$G\approx 20$ --- while leaving the precision quadrature, the evidence accumulation,
and the $W_g$ log-sum-exp untouched (they operate on the shared moments).

\paragraph{What is approximated.} The only $\tau$-dependence we drop is second order:
today, last step's $Q_g$ changes each node's analysis covariance, which changes the
resampled ensemble, which changes \emph{this} step's VSEM inputs. Because the wood-pool
process error $1/\tau$ is small relative to the state and VSEM is smooth, the
pushforward is nearly identical across nodes, so a shared ensemble plus per-node $Q_g$
is an excellent approximation. It is exactly the standard Rao--Blackwellized /
marginalized-particle construction: sample the state once, integrate the nuisance
parameter ($\tau$) analytically. If one is cautious, keep a handful of
\emph{representative} nodes with distinct ensembles (say the 3 highest-$W_g$) and share
the rest --- a tunable accuracy/speed knob.

\section{Idea 2 --- Replace the $J$-member forecast with a deterministic transform}
\label{sec:sigma}

Downstream, the forecast ensemble is used \emph{only} to form a mean and a $3\times3$
covariance (\code{colMeans(Nf[su,])}, \code{var(Nf[su,])},
\code{OR\_joint\_update.Rmd:346--347}); the members themselves are never propagated
(they are re-drawn from the collapsed analysis Gaussian each step). So there is no
reason to push $J\!\approx\!100$ random members through VSEM just to estimate two
moments.

\paragraph{Unscented transform.} For a 3-D state, the unscented transform uses
$2n+1=7$ deterministic sigma points and reproduces the pushforward mean and covariance
of a smooth map to second order. Replace the $J$-loop with 7 VSEM evaluations at the
sigma points, then reconstruct $(\hat\mu,\Sigma^{\text{ens}})$ with the standard UT
weights. This also \emph{removes the Monte-Carlo noise} that currently forces $J$ to be
large enough for a stable $3\times3$ \code{var()}.

\paragraph{Extended KF alternative.} Even cheaper: run VSEM once at the ensemble mean
and take a finite-difference Jacobian $F$ (3 extra runs), giving
$\Sigma^{\text{ens}} \approx F\,\Sigma_{\text{in}}\,F^\top$ --- about 4 evaluations.
Slightly less accurate than UT for strong nonlinearity; VSEM's annual pool dynamics are
mild, so either is defensible.

Combined with Idea~1, the per-step VSEM count drops from $G\cdot D\cdot J \approx
20\cdot5\cdot100 = 10^4$ to $D\cdot 7 = 35$.

\paragraph{Caveat.} The disturbance labels apply the stochastic operator
\code{disturbance.p} (\code{OR\_joint\_update.Rmd:353}) after VSEM. Its post-disturbance
biomass draw is itself Monte-Carlo; the UT/EKF pushforward should cover the pre-disturbance
VSEM step, with the disturbance-model moments added analytically (its mean/covariance are
available in closed form from $\mu_0$, $V_0$) rather than by resampling. This keeps the
disturbance branches cheap too.

\section{Idea 3 --- Prune negligible mixture components}
\label{sec:prune}

The mixture has $D\cdot G = 100$ components per node-set, but their weights are very
uneven, and forecasting a component whose weight is $\sim\!10^{-6}$ is wasted VSEM time.
The code already skips empty inputs (\code{if (bar.w[c, g] <= 0) next},
\code{OR\_joint\_update.Rmd:427}); extend this to an $\varepsilon$-threshold on the
effective weight.

\begin{itemize}[nosep]
  \item \textbf{Precision nodes.} After a few timesteps the node posterior $W_g$
        concentrates on a few $\tau$ values. Skip forecasting for nodes with
        $W_g < \varepsilon_W$ (they contribute negligibly to the final $\Pi_D$ and to
        $W_g\,\alpha^{c,k,g}$ at \code{OR\_joint\_update.Rmd:516}). Re-evaluate the
        active set each step so a node can re-enter if evidence shifts.
  \item \textbf{Disturbance labels.} The incoming weight of a disturbance branch is
        $\code{bar.w[c,g]}\cdot\beta_k$ with $\beta_k = p_t/\code{dist\$n}$
        (\code{OR\_joint\_update.Rmd:392}). In years with a tiny disturbance prior
        $p_t$, all four disturbance branches carry almost no mass; skip their VSEM +
        Kalman work when $\code{bar.w[c,g]}\cdot\beta_k < \varepsilon_\beta$.
  \item \textbf{Component merging.} Near-duplicate components across $c$ can be merged
        (moment-matched) \emph{before} the forecast, not just after, keeping
        $K_{\text{in}}$ small.
\end{itemize}
This is near-exact: with a conservative $\varepsilon$ the dropped mass is provably
below the threshold. \textbf{Log what was pruned each step} so a silent truncation never
masquerades as full coverage.

\section{Idea 4 --- Parallelize}
\label{sec:parallel}

Nothing in the cost model requires serial execution across the independent axes.
\begin{itemize}[nosep]
  \item \textbf{Sites} (\code{run\_joint\_update\_sites}, \code{OR\_joint\_update.Rmd:588})
        are embarrassingly parallel and each writes its own \code{.rds}. Swap the
        \code{for (site in site.idx)} loop for \code{parallel::mclapply} /
        \code{future.apply::future\_lapply}. This is the safest, largest win for batch
        runs and composes with everything else --- near-linear in cores.
  \item \textbf{The $g$-loop} (\code{OR\_joint\_update.Rmd:417}) is independent across
        nodes (evidence is combined only afterward), and \textbf{the $j$ VSEM loop}
        (\code{OR\_joint\_update.Rmd:333}) is independent across members. Either can be
        an \code{mclapply}. If Idea~1 is adopted the $g$-loop is no longer the VSEM
        bottleneck, so prefer parallelizing sites and/or the VSEM member evaluations.
\end{itemize}
Watch RNG: use \code{parallel}'s L'Ecuyer streams (\code{RNGkind("L'Ecuyer-CMRG")}) so
parallel runs are reproducible.

\section{Idea 5 --- Make VSEM itself cheaper}
\label{sec:vsem}

\code{BayesianTools::VSEM} is an R-level wrapper with per-call overhead, invoked $J$
times in a tight loop.
\begin{itemize}[nosep]
  \item \textbf{Batch the call.} VSEM's dynamics are a simple daily update; a vectorized
        or Rcpp reimplementation that accepts a \emph{matrix} of parameter sets and
        returns all members from one call eliminates both the per-call wrapper overhead
        and the R-level \code{for (j in 1:J)} loop. This is a constant-factor win
        ($2\text{--}10\times$) that stacks on the structural ideas. (Largely moot if
        Idea~2 shrinks $J$ to $\sim 7$, but still helps the sigma-point evaluations.)
  \item \textbf{Avoid redundant re-slicing.} \code{PAR[tsel]} and the
        \code{tsel <- (t-1)*365 + 1:365} index can be precomputed once per $t$; minor,
        but free.
\end{itemize}

\section{Idea 6 --- Micro-optimizations (exact, low-risk)}
\label{sec:micro}

\begin{itemize}[nosep]
  \item \textbf{Preallocate the accumulators.} \code{allk.*}, \code{allf.*},
        \code{up.*}, \code{fc.acc.*} grow with \code{c(...)} inside the loops
        (\code{OR\_joint\_update.Rmd:432--469}). Repeated \code{c()} reallocates and
        copies --- $\bigO(n^2)$ over a step. Preallocate fixed-size lists (length
        $D\cdot G\cdot K_{\text{in}}$) and fill by index.
  \item \textbf{Hoist invariants.} When $R$ is constant, precompute $R^{-1}$ and the
        $H$-projected quantities in \code{KalmanUpdate} (\code{OR\_joint\_update.Rmd:235})
        outside the loops instead of re-forming them per component.
  \item \textbf{Cheaper categorical draw.} Replace
        \code{apply(rmultinom(J,1,p),2,which.max)} (\code{OR\_joint\_update.Rmd:336})
        with a single \code{sample(0:dist\$n, J, replace=TRUE, prob=...)}.
  \item \textbf{Symmetrize once.} The \code{(Ck+t(Ck))/2 + diag(1e-6,3)} guard
        (\code{OR\_joint\_update.Rmd:491}) is fine; just avoid the \code{method="svd"}
        \code{rmvnorm} cost once $J$ is small (Idea~2).
\end{itemize}

\section{Recommended plan}
\label{sec:plan}

Order the work by payoff-to-risk:
\begin{enumerate}[nosep]
  \item \textbf{Parallelize sites (Idea~4).} Immediate, exact, near-linear; no change to
        results. Do this first for batch runs.
  \item \textbf{Prune nodes and labels (Idea~3).} Near-exact, easy, and $2\text{--}5\times$;
        log the dropped mass.
  \item \textbf{Rao--Blackwellize the precision (Idea~1).} The structural $\sim\!20\times$.
        Validate against the current filter on a few sites (compare $\Pi_D$, \code{probD},
        NIS calibration) before trusting it broadly.
  \item \textbf{Sigma-point forecast (Idea~2).} Another $\sim\!10\times$ and removes MC
        noise; validate the same way.
  \item \textbf{Batch/Rcpp VSEM and micro-opts (Ideas~5--6).} Constant-factor polish once
        the structure is settled.
\end{enumerate}

\paragraph{Validation.} Ideas 1, 2 and 3 change the numbers (1--2 are approximations,
3 drops sub-$\varepsilon$ mass), so before adopting them, run the current and the fast
version on the same handful of sites and check that (i) the posterior disturbance
probability \code{probD} and the overall analysis mean/CI agree within Monte-Carlo error,
and (ii) the NIS calibration (the $\chi^2_T/T$ band derived in
\code{claude-nis-6.13.26.tex}) is unchanged. Ideas 4, 5 and the micro-optimizations are
exact and need only a regression check that outputs are bit-for-bit (or RNG-stream)
identical.

\paragraph{Note on saved runs.} The existing \code{.rds} outputs in
\code{experiments/jointUpdateNewTest} predate the recent algorithm fixes (the
\code{bar.*} carry-forward, the full-mixture $\Pi_D$/NIS, and the analytic $q_{33}$).
Any speed refactor should regenerate them, and that regeneration is the natural moment
to fold in Ideas 1--4.

\end{document}
